\section{Integer sequences}

\begin{lemma} \label{lem:sequence argument1}
\lean{sequence_argument1}
\leanok
Let $n$ be a positive integer.
Let $m_1,\dots,m_n$ be a sequence of integers.
Suppose that there are no indices $1 \leq i < j \leq n$ such that
\begin{equation} \label{eq:sequence condition}
m_k = m_i + 1 \quad (i < k < j), \qquad m_j > m_i + 1.
\end{equation}
Then for $1 \leq i < j \leq n$,
\[
m_j \leq m_i+1 \qquad (1 \leq i < j \leq n).
\]
\end{lemma}
\begin{proof}
\leanok
We prove the claim by induction on $j-i$.
The base case $j-i=1$ is covered by the second part of \eqref{eq:sequence condition}.
For the induction step, suppose by way of contradiction that $m_j> m_i+1$ for some $i,j$.
Set $m := m_i + 1$.
For $k = i+1, \dots, j-1$ we have $m_{k} \leq m_i+1$ and $m_j \leq m_{k}+1$ by the induction hypothesis;
combining with $m_j \geq m_i + 2$ shows that both of these must be equalities
and so $m_k = m, m_j = m+1$.
Now $(m_i, \dots, m_j) = (m-1, m, \dots, m, m+1)$, which contradicts \eqref{eq:sequence condition}.
\end{proof}

\begin{lemma} \label{lem:sequence argument2}
\lean{sequence_argument2}
\leanok
Let $n$ be a positive integer. Let $c$ be an integer.
Let $m_1,\dots,m_n$ be a sequence of integers satisfying the following conditions.
\begin{enumerate}
\item[(a)]
We have $m_1 + \cdots + m_n = c$.
\item[(b)]
For $1 \leq i < j \leq n$, $m_j \leq m_i+1$.
\end{enumerate}
Then $m_n \leq \lfloor \frac{c+n-1}{n} \rceil$ (i.e., $\lceil \frac{c}{n} \rceil$).
\end{lemma}
\begin{proof}
\leanok
For $n = 1$, note that (a) directly forces $m_1 = 0$.
For $n > 1$, we have $m_n = (m_1 + \cdots + m_n) - (m_1 + \cdots + m_{n-1}) \geq 0$ by (a) and (b).
Moreover, if $m_n \geq 1$, then by (c) we have $m_1,\dots,m_{n-1} \geq 0$ and hence
$m_1 + \cdots + m_n \geq 1$, a contradiction. Hence $m_n \leq 0$ as claimed.
By (b), we have 
\[
c = m_1 + \cdots + m_n \geq (n-1)(m_n - 1) + m_n = nm_n - n + 1.
\]
We thus have $m_n \leq \lfloor \frac{c+n-1}{n}\rfloor$.

\end{proof}

\begin{lemma} \label{lem:sequence argument3}
\lean{sequence_argument3}
\leanok
Let $n$ be a nonnegative integer.
Let $c$ be an integer.
Let $S$ be a subset of $\Z^n$ satisfying the following conditions.
\begin{enumerate}
\item[(a)]
The set $S$ is nonempty.
\item[(b)]
For every tuple $(m_1,\dots,m_n) \in S$, we have $m_1 + \cdots + m_n = c$.
\item[(c)]
For $i=1,\dots,n-1$, there exists a constant $c_i \in \Z$ such that
for every tuple $(m_1,\dots,m_n) \in S$, $m_1 + \cdots + m_i \leq c_i$.
\item[(d)]
For every tuple $(m_1,\dots,m_n) \in S$,
if there exist indices $1 \leq i < j \leq n$ such that
\begin{equation} \label{eq:sequence condition3}
m_k = m_i + 1 \quad (i < k < j), \qquad m_j \geq m_i + 2,
\end{equation}
then there exists another tuple $(m'_1,\dots,m'_n) \in S$ with
\[
m'_1 + \cdots + m'_i > m_1 + \cdots + m_i, \qquad
m'_1 + \cdots + m'_k \geq m_1 + \cdots + m_k \quad (k \neq i).
\]
\end{enumerate}
Then there exists a tuple $(m_1,\dots,m_n) \in S$ with $m_n \leq \lfloor \frac{c+n-1}{n} \rceil$.
\end{lemma}
\begin{proof}
\leanok
\uses{lem:sequence argument1,lem:sequence argument2}
Define the function $f\colon S \to \Z$ by
\[
f(m_1,\dots,m_n) = \sum_{i=1}^{n-1} \sum_{j=1}^i m_j;
\]
from conditions (b) and (c), we see that $\sum_{j=1}^i m_i \leq c_i$ for each $i$, so $f$ is bounded above by $\sum_{i=1}^{n-1} c_i$.
(We also have the simpler expression $f(m_1,\dots,m_n) = \sum_{i=1}^n (n-i) m_i$ but we will not need it.)

Since $S$ is nonempty, we can choose a tuple $(m_1,\dots,m_n) \in S$ on which $f$ achieves its maximum.
There then cannot exist indices $i,j$ satisfying \eqref{eq:sequence condition3}, as then the new sequence $(m'_1,\dots,m'_n)$ would have $\sum_{i=1}^{n-1} (n-i)m'_i > \sum_{i=1}^{n-1} (n-i)m_i$.
Consequently, we can apply Lemma~\ref{lem:sequence argument1} to deduce that 
$m_j \leq m_i + 1$ for $1 \leq i < j \leq n$, then Lemma~\ref{lem:sequence argument2} to conclude.
\end{proof}

